% 第9章 光子拡散と減衰
\section{光子拡散と減衰}
\label{ch:ch09_damping}

\paragraph{この章で導くこと（入力→出力）}
強結合展開の次数を上げ（入力）、せん断粘性と熱伝導による拡散減衰（Silk damping）を
導いて減衰スケール $k_D(\eta)$ を定量化する（出力 S9.1）。これが第\ref{ch:ch12_peaks}章の
減衰尾を説明する。

\subsection{平均自由行程とランダムウォークの直観}
光子は電子と衝突しながらランダムウォークで拡散する。共動拡散長は
$\lambda_D^2\sim\int\lambda_{\rm mfp}\,c\,dt/a$ のオーダーで、平均自由行程
$\lambda_{\rm mfp}=1/(n_e\sigma_T a)$ が再結合に向けて伸びるほど拡散が進む。
スケール $k>k_D\sim1/\lambda_D$ のゆらぎは散乱でならされて消える。

\subsection{強結合1次の系統的展開（導出契約 S9.1）}
$1/\tau'$ 展開の1次で、四重極 $\Theta_2$（粘性）と $\Theta_1+v_b/3$ のずれ（滑り＝
熱伝導）が現れる。これを振動子方程式（第\ref{ch:ch08_acoustic}章）に戻すと分散関係に
虚部が生じ、振幅が $\exp[-(k/k_D)^2]$ で減衰する。減衰スケールは
\begin{equation}
  \frac{1}{k_D^2(\eta)}=\int_0^\eta\frac{d\eta'}{6(1+R)\,|\dot\tau|}
  \left[\frac{R^2}{1+R}+\frac{16}{15}\right],
  \label{eq:kD}
\end{equation}
で与えられる。\textbf{偏光を無視すると最後の係数が $16/15$ になる}（偏光込みでは
$8/9$）。この差が CLASS 比での減衰尾の系統差（第\ref{ch:ch14_numerics}章・付録E）の
主因である。

\begin{numreality}
式\eqref{eq:kD} は \texttt{cmbcore.analytic.silk\_scale} に実装。fiducial で
\begin{equation}
  k_D(\eta_*)\approx0.149~\mathrm{Mpc^{-1}}\ (\text{目標}\ 0.14\pm20\%),
\end{equation}
を与え、対応する減衰多重極は $\ell_D\sim k_D\chi_*\sim1300$--$1500$。
\end{numreality}

\subsection{可視度幅によるぼかし}
最終散乱面が有限の厚み $\Delta z\approx80$（第\ref{ch:ch03_recombination}章）を持つため、
視線方向に異なる位相の振動が重なり、さらに小スケールがならされる（Landau 減衰的）。
これも $e^{-k^2/k_D^2}$ 型の抑制に寄与する。

\subsection{数値確認}
$k=10^4 H_0/c$ のような小スケールでは $\Theta_0$ が急減衰する（NB3）。減衰包絡
$e^{-(k/k_D)^2}$ とフル数値の振幅は $k\lesssim0.2$~Mpc$^{-1}$ で $20\%$ 以内に一致する
（NB7, F7.3）。

\paragraph{まとめ}
拡散減衰スケール $k_D$\eqref{eq:kD} を導き、$k_D\approx0.149$~Mpc$^{-1}$ を数値確認した。
偏光無視による係数 $16/15$ が減衰尾の既知系統を生む。

\paragraph{演習}
(1) 式\eqref{eq:kD} の次元を確認せよ。
(2) NB7 で $e^{-(k/k_D)^2}$ と数値振幅を重ね、ずれが出る $k$ を同定せよ。
(3) $\ell_D=k_D\chi_*$ を数値で求め減衰尾の落ち始めと比べよ。
